\makeabstract
    {
    Investigating Novel Synthesis of Organic Radical Multiferroics
    }
    {
    Emmett Murphy\textsuperscript{1}, Zen Loo\textsuperscript{1}, Ren A. Wiscons\textsuperscript{1}
    }
    {
    \textsuperscript{1} Department of Chemistry, Amherst College, Amherst, MA
    }
    {
    Previous research in the Wiscons Lab has explored the peculiar characteristics of ''triangulenes,'' or heteroatom-centered, triply oxygen-bridged triphenyl structures, seeking to work towards organic data storage. This study seeks a new route of synthesis for a carbon radical-centered triangulene that, due to its unpaired electron, should possess an electric and magnetic dipole. A material with both of these qualities is known as ''multiferroic,'' and organic multiferroics are promising candidates for highly manipulable materials.
    }
    {
    To synthesize precursor 1,3-Bis-(3-methoxyphenoxy)benzene, we attempted a copper-catalyzed Ullmann coupling reaction to combine reagents resorcinol and 3-iodoanisole. This reaction ran overnight for 20 hours at 135{\textdegree}C. After reacting, the product was first pulled through a celite plug, then washed with 0.5M HCl in a separatory funnel. The organic layer was then dried with magnesium sulfate, and the product extracted via column chromatography.
    }
    {
    Though the scale of the reaction is small, we achieved a relatively high yield; 1H NMR testing confirmed the product's identity and high degree of purity. 3.1 grams of precursor were successfully synthesized over 8 reactions.
    }
    {
    We have accrued enough precursor to proceed to the next steps: lithiation and ring-closing. Because we cannot safely work with n-butyllithium ourselves, Dr. Brian H. Northrop and his group at Wesleyan University have graciously agreed to perform the lithiation and nucleophilic ring-closing steps. We hope that this will yield a triangulene cation that can be subsequently reduced to a radical in future work.
    }

    \newpage
    